\begin{textsample}{POS Dim 1 – ms – Score 58.00 – ms/ms\_em\_29.txt}  \label{ex:f1_pos_183}
3.1 Área de Conhecimento : Linguagens e suas Tecnologias A Base Nacional Comum Curricular ( BNCC ) determina que a área de Linguagens e suas Tecnologias explore os diversos tipos de linguagem, expandindo as \textbf{capacidades} de expressões artísticas, corporais e linguísticas dos estudantes.
A participação dos \textbf{discentes} nas diferentes práticas de linguagem colabora para sua formação integral, que não se deve limitar ao aperfeiçoamento de \textbf{competências} cognitivas, mas também das \textbf{socioemocionais}, a fim de que a aprendizagem \textbf{auxilie} na resolução de \textbf{problemas} e produção de conhecimentos que reflitam no desenvolvimento da sociedade.
A educação integral é um dos princípios específicos da etapa do Ensino Médio e todas as suas modalidades, conforme determina o artigo 5º, inciso I da Resolução n. 3, de 21 de novembro de 2018, que atualiza as Diretrizes Nacionais para o Ensino Médio ( DCNEM ).
Assim, o que se pretende é que a escola prepare o estudante para aprender continuamente, aplicar os conhecimentos construídos em seu benefício e de uma coletividade.
No artigo 6, inciso I, define -se : > [... ] formação integral : é o desenvolvimento intencional dos aspectos físicos, cognitivos e \textbf{socioemocionais} do estudante por meio de processos educativos significativos que promovam a autonomia, o comportamento cidadão e o protagonismo na construção de seu projeto de vida ( BRASIL, 2018a ).
Segundo a Base Nacional Comum Curricular, ao desenvolver as \textbf{competências} e habilidades da área de Linguagens e suas Tecnologias, o estudante tem a oportunidade de consolidar e ampliar sua \textbf{capacidade} “ de uso e reflexão sobre as linguagens – artísticas, corporais e verbais ” ( BRASIL, 2018c, p. 481 ), ou seja, não será somente leitor, mas \textbf{autor} de diversos gêneros \textbf{discursivos} \textbf{contemporâneos}, muitas vezes, relacionados ao universo \textbf{digital}.
Cabe ressaltar que a área de Linguagens e suas Tecnologias ( LGG ) propõe que se ultrapasse a linguagem verbal, considerando todos os elementos que contribuem para a significação do texto.
A proposta é de uma escola que \textbf{aproxime} os estudantes das mais diversas produções, em especial, as artísticas e culturais, dentro dos contextos regionais.
Conforme afirma Rojo : > Assim como foi capaz de popularizar os impressos, urge que a escola se \textbf{preocupe} com o acesso a outros espaços valorizados de cultura ( museus, bibliotecas, teatros, espetáculos ) e a outras mídias ( analógicas e \textbf{digitais} ) ( ROJO, 2009, p. 52 ).
Na etapa do Ensino Médio, a BNCC define como prioridade cinco campos de atuação social, que orientam para práticas de linguagem que possibilitem aos estudantes “ vivenciarem experiências significativas em diferentes mídias ( impressa, \textbf{digital}, analógica ) ”.
Os campos de atuação objetivam o avanço da educação escolar ao propor que a cultura \textbf{digital}, os \textbf{multiletramentos} e os novos \textbf{letramentos} sejam acrescidos à cultura do impresso ( escrita ), propiciem a \textbf{análise} de elementos visuais ( imagens estáticas ou em movimento ), sonoros ( músicas, ruídos, sonoridades ), verbais ( oral ou \textbf{visual-motora}, como \textbf{Libras} e escrita ) e corporais ( gestuais, cênicas, dança ), ou seja, considerem diferentes semioses.
Esses campos de atuação compreendem : - campo da vida pessoal : processos de construção de identidade e de Projetos de Vida ; - campo das práticas de estudo e pesquisa : construção do conhecimento científico para aprender a aprender ; - campo jornalístico-midiático : construção de consciência \textbf{crítica} e seletiva em relação à produção e circulação de informações ; - campo de atuação na vida pública : reflexão e participação na vida pública, pautado pela ética ; - campo artístico : reconhecimento, valorização, fruição e produção de manifestações artísticas em geral.
Nessa perspectiva, o Currículo de Referência de Mato Grosso do Sul — Ensino Médio da área de Linguagens e suas Tecnologias — estrutura -se em um Organizador Curricular, que corresponde a um \textbf{quadro} composto por cinco ( 05 ) colunas, contendo : eixo temático ( que em LGG equivale ao campo de atuação social ), habilidades, componente curricular, objeto de conhecimento e sugestões didáticas, os quais apresentam os saberes que devem ser aprofundados com os estudantes, por meio de \textbf{competências} e habilidades, em cada componente curricular da Área de Conhecimento, de determinado ano, dando respaldo pedagógico aos professores.
Considerando que os estudantes possuem ritmos de aprendizagem muito diferentes uns dos outros, as graduações das complexidades das habilidades devem acompanhar o desenvolvimento de cada indivíduo.
O Estado de Mato Grosso do Sul optou por manter os componentes curriculares, já existentes, na formação geral, partindo de sugestões didáticas que desenvolvam as \textbf{competências} e habilidades \textbf{tanto} por Área de Conhecimento quanto interdisciplinarmente.
No espaço do \textbf{organizador} curricular da Formação Geral Básica ( FGB ), destinado às sugestões didáticas, são apresentadas práticas \textbf{metodológicas} que proporcionam aos professores orientações de estudo por área.
Enquanto sugestões, é importante destacar que essas ações não devem ser compreendidas como obrigatórias, mas como possibilidades que precisam fazer sentido para as unidades escolares, que podem adaptá -las ou usá -las como inspiração, ou definir suas próprias ações didáticas, considerando as especificidades de cada \textbf{localidade}.
Essas sugestões estão atreladas aos Temas \textbf{Contemporâneos}, à prática da presença pedagógica, ao uso de metodologias ativas, à pesquisa científica, ao protagonismo juvenil e à \textbf{autoria}, assegurando, assim, o desenvolvimento global do estudante e a construção de seu Projeto de Vida, conforme previsto no artigo 11 das DCNEM : > Art. 11 A formação geral básica é composta por \textbf{competências} e habilidades previstas na Base Nacional Comum Curricular ( BNCC ) e articuladas como um todo indissociável, enriquecidas pelo contexto histórico, \textbf{econômico}, social, ambiental, cultural local, do mundo do trabalho e da prática social, e deverá ser organizada por áreas de conhecimento [... ] ( BRASIL, 2018a ).
Assim, as atividades \textbf{sugeridas} em LGG \textbf{mobilizam} conjuntos de \textbf{competências} e habilidades que se \textbf{inter-relacionam}, denominados agrupamentos, possibilitando o \textbf{aprofundamento} do conhecimento, a partir de estratégias \textbf{metodológicas} que se complementam com os componentes curriculares da área com as demais, quando possível.
Destaca -se que cada componente ainda explora suas especificidades, proporcionando ao estudante a \textbf{capacidade} de identificar, analisar e posicionar -se, criticamente, diante de situações da realidade.
No Estado de Mato Grosso do Sul, a área de Linguagens e suas Tecnologias é composta pelos seguintes componentes curriculares : Arte, Educação Física, Língua \textbf{Inglesa} e Língua Portuguesa, na FGB, e Língua \textbf{Espanhola} no Núcleo Integrador¹⁰, para casos em que a escola ofereça uma segunda língua estrangeira, conforme o que define a Lei 13.415, de 2017, em seu artigo 35 A, \textbf{parágrafo} 4º : > Os currículos do ensino \textbf{médio} incluirão, obrigatoriamente, o estudo da língua \textbf{inglesa} e poderão ofertar outras línguas estrangeiras, em caráter optativo, preferencialmente o \textbf{espanhol}, de acordo com a disponibilidade de oferta, locais e horários definidos pelos sistemas de ensino ( BRASIL, 2017 ).
Ressalta -se que essa organização tem como objetivos potencializar a construção do conhecimento e articular os componentes curriculares, partindo de um ponto em comum, nesse caso, das \textbf{competências} e habilidades a serem desenvolvidas, bem como do campo de atuação social.
Para a área de Linguagens e suas Tecnologias apresentam -se dois \textbf{organizadores} curriculares : o primeiro explora o desenvolvimento de habilidades e \textbf{competências}, por área, e o segundo enfatiza as habilidades específicas de Língua Portuguesa, voltadas às práticas de leitura, escrita e oralidade, de diversos gêneros \textbf{discursivos} e ao estudo do texto literário, tendo em \textbf{vista} os mecanismos linguísticos e \textbf{multissemióticos}, visando à formação do “ lautor ” ( ROJO, 2013, p. 20 ). ¹⁰ O Núcleo Integrador compreende os componentes curriculares Projeto de Vida, Empreendedorismo Social, Ciências Integradas e Novas Tecnologias, Matemática Criativa e Resolução de \textbf{Problemas}, Linguagens e Interartes, Língua \textbf{Espanhola} e Intervenção Comunitária que se articulam de forma integrada e \textbf{permeiam} por todas as áreas de conhecimento e possibilitam o desenvolvimento de \textbf{competências} e habilidades específicas da BNCC e dos Itinerários Formativos. ( MATO GROSSO DO SUL, 2020 ).
O desenvolvimento das \textbf{competências} e habilidades foi definido agrupando aquelas que tinham uma relação próxima, de forma que pudessem ser trabalhadas em uma \textbf{sequência} didática que passasse pelos principais níveis de aprendizagem, conforme a Taxonomia de Bloom revisada em 2001, chegando ao sexto nível : o da criação e produção ( \textbf{figura} 1 ). \textbf{Figura} 1 - Taxonomia de Bloom revisada O processo de \textbf{seriação} foi realizado da seguinte \textbf{maneira} : agrupamento de habilidades que se completam ; definição de objetos de conhecimento ; produção de sugestões didáticas ; \textbf{análise} do grau de dificuldade dos agrupamentos de habilidades por componente ; e alinhamento por área do conhecimento e \textbf{seriação}.
Nesse sentido, todas as sugestões didáticas oferecem ao menos uma forma de materialização do conhecimento, a partir de uma produção que utilize as linguagens verbal, corporal, visual e/ou sonora explorando, assim, diferentes semioses e empregando, em diversas ocasiões, as ferramentas próprias da cultura \textbf{digital}.
No \textbf{organizador}, os componentes curriculares aparecem por ordem alfabética : Arte, Educação Física, Língua \textbf{Espanhola}, Língua \textbf{Inglesa} e Língua Portuguesa.
No componente curricular Arte, etapa do Ensino Médio, conforme a BNCC, o arte educador promove : > [... ] o entrelaçamento de culturas e saberes, possibilitando aos estudantes o acesso e a interação com as distintas manifestações culturais populares presentes na sua comunidade.
O mesmo deve ocorrer com outras manifestações presentes nos centros culturais, museus e outros espaços, de modo a propiciar o exercício da \textbf{crítica}, da apreciação e da fruição de \textbf{exposições}, concertos, apresentações musicais e de dança, filmes, peças de teatro, poemas e obras literárias, entre outros, garantindo o respeito e a valorização das diversas culturas presentes na formação da sociedade brasileira, especialmente as de \textbf{matrizes} indígena e africana ( Brasil, 2018c, p. 483 ).
Nesse contexto, o Currículo considera as peculiaridades da comunidade escolar, prioriza a valorização de produções locais, por meio da arte regional, nacional e mundial, para que ocorra a identificação cultural e a visão \textbf{crítica} de sua existência, respeitando a diversidade das culturas locais do Estado de Mato Grosso do Sul.
O objetivo da Arte na Educação Básica não é formar artistas, mas garantir que conhecimentos produzidos contribuam para a \textbf{contextualização} de saberes e práticas artísticas do estudante, valorizando a arte local, a produção cultural e a construção de seu Projeto de Vida.
Ressalta -se a necessidade de um olhar flexível, autônomo e coerente, \textbf{focado} no protagonismo e na articulação das práticas de : criar, ler, produzir, construir, exteriorizar e refletir sobre formas artísticas, comprometido com “ o \textbf{aprofundamento} na pesquisa e no desenvolvimento de processos de criações autorais nas linguagens das artes visuais, do audiovisual, da dança, do teatro, das artes circenses e da música ” ( BRASIL, 2018c, p. 482 ), conforme Barbosa : > [... ] Através das artes é possível desenvolver a percepção e a imaginação, apreender a realidade do meio ambiente, desenvolver a \textbf{capacidade} \textbf{crítica}, permitindo analisar a realidade percebida e desenvolver a criatividade de \textbf{maneira} a mudar a realidade que foi analisada ( BARBOSA, 1998. p. 16 ).
Para o componente curricular Educação Física, na etapa do Ensino Médio, pretende -se que o estudante aprofunde as experiências do Ensino Fundamental sobre as \textbf{capacidades} e limites do corpo ; compreendendo a importância do estilo de vida saudável e do autocuidado na manutenção da saúde.
Dessa forma, o processo vivenciado será consolidado na última fase da Educação Básica : > [... ] ao \textbf{final} do Ensino Médio, o jovem deverá apresentar uma \textbf{compreensão} aprofundada e sistemática acerca da presença das práticas corporais em sua vida e na sociedade, incluindo os fatores sociais, culturais, ideológicos, \textbf{econômicos} e políticos envolvidos nas práticas e nos discursos que circulam sobre elas ( BRASIL, 2018c, p. 495 ).
Ao pensar nas \textbf{competências} e habilidades relacionadas à aprendizagem da área de Linguagem e suas Tecnologias, no que tange às particularidades do Componente Curricular Educação Física, logo tem -se em \textbf{mente} os movimentos corporais diversos : esportes, caminhadas, corridas, treinamentos e outros, no entanto é possível associá -los ao uso das Tecnologias \textbf{Digitais} da Informação e Comunicação ( TDIC ) ao propor novas estratégias \textbf{metodológicas}.
Os videogames, geralmente tidos como inimigo da vida saudável, podem \textbf{auxiliar} na manutenção da qualidade de vida a partir de jogos que favoreçam a atividade física ( exergames ).
Segundo Silva : > Atualmente os Exergames estão sendo comparados a boas atividades de movimentos e práticas, atividade de estímulo ao prazer de se praticar uma atividade e ao mesmo tempo de possibilidade de aprimoramento do aprendizado de crianças e adolescentes em escolas ( SILVA, 2017, p. 6 ).
Assim, propõem -se atividades que \textbf{auxiliam} o professor a compor uma \textbf{sequência} didática que colabore com a qualidade de vida, sem ignorar os avanços tecnológicos ao desenvolver as \textbf{competências} e habilidades da área.
No que \textbf{diz} respeito à Língua \textbf{Inglesa}, a Lei de Diretrizes e Bases, redação dada pela Lei n. 13.415/2017, tornou obrigatório o seu estudo, na etapa do Ensino Médio, \textbf{devido} ao seu “ caráter global – pela multiplicidade e variedade de usos, usuários e funções na \textbf{contemporaneidade} ” ( BRASIL, 2018c. p. 484 ).
Ao propor a aprendizagem de Língua \textbf{Inglesa}, por área de conhecimento, objetiva -se que os estudantes explorem as diversas situações de uso da língua, especialmente na cultura \textbf{digital}, ampliando sua \textbf{capacidade} \textbf{discursiva} e de reflexão, a partir da \textbf{contextualização} das práticas de linguagem nos diversos campos de atuação : > Trata -se, portanto, de expandir os repertórios linguísticos, \textbf{multissemióticos} e culturais dos estudantes, possibilitando o desenvolvimento de maior consciência e reflexão \textbf{críticas} das funções e usos do \textbf{inglês} na sociedade \textbf{contemporânea} – permitindo, por exemplo, problematizar com maior criticidade os motivos pelos quais ela se tornou uma língua de uso global ” ( BRASIL, 2018c.p. 485 ).
O componente curricular Língua Portuguesa, em consonância com a Lei 13.415, é obrigatório nos três anos da etapa do Ensino Médio, juntamente com a Matemática.
Pretende -se, nessa fase, aprofundar os conhecimentos adquiridos na Etapa do Ensino Fundamental referentes a linguagem ( uso, finalidade, significação ) e à leitura do texto literário.
Nesse sentido, a Língua Portuguesa, indissociável da Literatura, busca relacionar -se com as mais diversas linguagens, presentes em diferentes gêneros \textbf{discursivos}, a fim de que o estudante tenha condições de participar, significativamente, das práticas sociais \textbf{contemporâneas}.
Pensando nos níveis de alfabetismo definidos pelo Indicador Nacional de Alfabetismo \textbf{Funcional} — INAF, fica claro que é necessário, ao concluir a Etapa do Ensino Médio, que o estudante tenha “ \textbf{capacidade} de ler textos longos, orientando -se por subtítulo, localizando mais de uma informação, de acordo com as condições estabelecidas, relacionando partes de um texto, comparando dois textos, realizando inferências e sínteses ” ( ROJO, 2009, p.47 ).
Ao construir as sugestões didáticas em Língua Portuguesa, considerou -se o desenvolvimento das práticas de linguagens que, na etapa do Ensino Médio, são chamadas de eixos de integração — leitura, produção de texto, oralidade ( escuta e produção oral ) e \textbf{análise} linguística/semiótica.
Faz -se necessário \textbf{lembrar} que, em alguns momentos, um ou mais eixos de integração estará/estarão em maior evidência, \textbf{devido} ao objeto de conhecimento que será tratado, ou seja, ora enfatiza -se a leitura, ora a oralidade, ora a produção textual.
Com o objetivo de continuar a formação do leitor e ampliar o contato dos estudantes com a Literatura, a Língua Portuguesa explora a riqueza proporcionada por \textbf{análises} contextualizadas, principalmente das obras clássicas, com momentos de investigação, \textbf{correlação} de saberes, resultando no pensamento \textbf{crítico} sobre a realidade.
O estudo da Literatura deve \textbf{supor} o entendimento da linguagem literária enquanto construção linguística que tem sua especificidade histórica, cultural e artística, e deve ser entendido e articulado com as formas de conhecimento e com as condições históricas de sua produção.
Assim, o texto literário é compreendido como \textbf{instância} que reúne uma realidade estética, uma forma de conhecimento que tem como finalidade entender o ser, a existência, o mundo e sua produção de vida nas variadas dimensões de figuração da vida social.
É importante compreender que a literatura faz parte da formação humana e sempre esteve presente no imaginário de cada pessoa.
Segundo Cândido, quando se trata de literatura : > Não há povo e não há homem que possa viver sem ela, isto é, sem a possibilidade de entrar em contacto com alguma espécie de fabulação.
Assim como todos sonham todas as noites, ninguém é capaz de passar as vinte e quatro horas do dia sem alguns momentos de entrega ao universo fabulado ( CÂNDIDO, 1995. p.174 ).
Já em relação ao componente Língua \textbf{Espanhola}, esse encontra -se próximo à realidade cultural e \textbf{econômica} do Estado de Mato Grosso do Sul, \textbf{devido} a sua localização geográfica e suas fronteiras com a Bolívia e o Paraguai.
Para as escolas que optarem pela oferta de uma segunda língua estrangeira, este Currículo de Referência de Mato Grosso do Sul propõe estratégias \textbf{metodológicas} alinhadas aos demais componentes curriculares da área de Linguagens e suas Tecnologias, o que viabiliza uma ação integrada e valoriza essa língua que \textbf{influencia}, fortemente, a cultura regional.
A área de Linguagens e suas Tecnologias contempla 7 ( sete ) \textbf{competências} específicas e 28 ( vinte e oito ) habilidades, apresentando, ainda, 54 ( cinquenta e quatro ) habilidades de Língua Portuguesa, consideradas \textbf{desdobramentos}, o que possibilita o planejamento por área do conhecimento.
Ressalta -se que foram criadas mais duas habilidades, uma para Linguagens e suas Tecnologias ( MS.EM13LGG6.n.01 ) e outra para Língua Portuguesa ( MS.EM13LP4.n.01 ).
As \textbf{competências} e habilidades estão organizadas considerando a proximidade existente entre elas, de formas que possam ser abordadas em uma situação didática que contemple os principais níveis de aprendizagem, conforme a Taxonomia de Bloom, chegando ao sexto nível : o da criação, produção.
O processo de \textbf{seriação} foi realizado da seguinte \textbf{maneira} : 1. agrupamento de habilidades que se completam ; 2. definição de objetos de conhecimento ; 3. produção de sugestões didáticas ; 4. \textbf{análise} do grau de dificuldade dos agrupamentos de habilidades por componente ; 5. alinhamento por área do conhecimento e \textbf{seriação}.
Dessa forma, todas as sugestões didáticas oferecem ao menos uma forma de materialização do conhecimento, a partir de uma produção que explore as linguagens verbal, corporal, visual e/ou sonora, e diferentes semioses, com o emprego das ferramentas próprias da cultura \textbf{digital}.

% matched lemmas: análise, aprofundamento, aproximar, autor, autoria, auxiliar, capacidade, competência, compreensão, contemporaneidade, contemporâneo, contextualização, correlação, crítico, desdobramento, devido, digital, discente, discursivo, dizer, econômico, espanhol, exposição, figura, figurar, final, focar, funcional, influenciar, inglês, instância, inter-relacionar, lembrar, letramento, libra, localidade, maneira, matriz, mente, metodológico, mobilizar, multiletramento, multissemiótico, médio, organizador, parágrafo, permear, preocupar, problema, quadro, sequência, seriação, socioemocional, sugerir, supor, tanto, vista, visual-motor
\end{textsample}
